% !TeX root = ../../main.tex
\section{压缩与还原流形拟合}\label{cw:sec:ml-compress}

当观测取值于高维向量空间 $\mathcal S_X=\mathbb R^{n_D}$ 时，直接在环境空间中实施似然加权或离散耦合输运会遭遇严重的维数灾难。实际物理或视觉信号往往并不充满整个环境空间，而是集中于某个自由度仅为 $n_d$（$n_d<n_D$）的低维子集 $\mathcal M\subseteq\mathcal S_X$ 附近。本节从经验模内积下的重构误差出发，给出压缩算子 $\mathscr T_{\Rightarrow}:\mathcal S_X\to\mathcal S_Z$（$\mathcal S_Z=\mathbb R^{n_d}$）与还原算子 $\mathscr T_{\Leftarrow}:\mathcal S_Z\to\mathcal S_X$ 的流形拟合框架，为第~\ref{cw:sec:algorithms}~章的潜空间输运滤波提供低维状态坐标与经验推前先验。

设离线观测样本矩阵为 $x_{1..M}=(x_1,\dots,x_M)\in\mathbb R^{n_D\times M}$。在样本行排列空间上定义经验模内积 $\langle\bm U,\bm V\rangle_{\hat{\mathbb P}_M}:=\frac1M\bm U\bm V^{\mathrm T}\in\mathbb R^{n_D\times n_D}$ 及其诱导范数 $\|\bm U\|_{\hat{\mathbb P}_M}^2:=\operatorname{tr}\langle\bm U,\bm U\rangle_{\hat{\mathbb P}_M}$。对任意重构映射 $\mathscr R:\mathcal S_X\to\mathcal M$，记余映射为 $\mathscr R^{\perp}:=\mathscr I-\mathscr R$，则样本平均重构平方误差可统一表示为
\begin{equation}\label{cw:eq:manifold-emp-norm}
\frac1M\sum_{i=1}^M\bigl\|x_i-\mathscr R(x_i)\bigr\|^2
=\operatorname{tr}\bigl\langle\mathscr R^{\perp}(x_{1..M}),\mathscr R^{\perp}(x_{1..M})\bigr\rangle_{\hat{\mathbb P}_M}
=\bigl\|\mathscr R^{\perp}(x_{1..M})\bigr\|_{\hat{\mathbb P}_M}^2.
\end{equation}
若不对 $\mathscr R$ 施加维数瓶颈而直接极小化式~\eqref{cw:eq:manifold-emp-norm}，最优解将平凡退化为恒等算子 $\mathscr I$。引入以 $\mathcal S_Z$ 为中介的压缩—还原复合结构 $\mathscr R=\mathscr T_{\Leftarrow}\circ\mathscr T_{\Rightarrow}$，即可将重构像集限制在至多 $n_d$ 个自由度的低维几何结构上。

\subsection{平直线性子空间的伴随投影与谱截断}

先考察 $\mathcal M$ 为 $n_d$ 维平直真子空间 $W<\mathcal S_X$ 的基础情形。此时压缩与还原算子均为线性映射，标准基下的表示矩阵分别记为 $\bm T_{\Rightarrow}\in\mathbb R^{n_d\times n_D}$ 与满列秩矩阵 $\bm T_{\Leftarrow}\in\mathbb R^{n_D\times n_d}$，像空间 $W:=\operatorname{im}(\mathscr T_{\Leftarrow})<\mathcal S_X$ 自然满足 $\operatorname{rank}(\bm T_{\Leftarrow}\bm T_{\Rightarrow})\le n_d$。记向 $W$ 的正交投影算子为 $\mathscr P_W$（矩阵表示为 $\bm P_W$），余投影为 $\mathscr P_W^{\perp}=\mathscr I-\mathscr P_W$（矩阵表示为 $\bm P_W^{\perp}=\bm I_{n_D}-\bm P_W$）。利用 $\operatorname{im}(\bm P_W^{\perp})\perp W$ 及迹的轮换性，重构误差可作正交分解：
\begin{equation}\label{cw:eq:manifold-orth-decomp}
\bigl\|(\bm I_{n_D}-\bm T_{\Leftarrow}\bm T_{\Rightarrow})x_{1..M}\bigr\|_{\hat{\mathbb P}_M}^2
=\operatorname{tr}\bigl(\bm P_W^{\perp}\langle x_{1..M},x_{1..M}\rangle_{\hat{\mathbb P}_M}\bigr)
+\bigl\|(\bm P_W-\bm T_{\Leftarrow}\bm T_{\Rightarrow})x_{1..M}\bigr\|_{\hat{\mathbb P}_M}^2.
\end{equation}
式~\eqref{cw:eq:manifold-orth-decomp} 右端首项仅依赖像子空间 $W$ 的选取，次项则度量复合映射偏离正交投影的程度。由此可直接导出线性压缩与子空间选取的闭式最优解，其完整推导见附录~\ref{cw:app:manifold}。

\begin{proposition}[线性压缩还原的伴随闭式解与谱截断]\label{cw:prop:linear-manifold}
设 $\mathscr T_{\Leftarrow}:\mathcal S_Z\to\mathcal S_X$ 为单射线性算子，$W=\operatorname{im}(\mathscr T_{\Leftarrow})$，$\mathscr T_{\Leftarrow}^{\star}:\mathcal S_X\to\mathcal S_Z$ 为其内积伴随算子。记样本二阶矩矩阵 $\langle x_{1..M},x_{1..M}\rangle_{\hat{\mathbb P}_M}$ 的特征值按降序排列为 $\lambda_1\ge\dots\ge\lambda_{n_D}\ge0$，相应规范正交特征向量族为 $u_{1..n_D}$。则：
\begin{enumerate}[label=(\arabic*),leftmargin=2.2em,topsep=2pt,itemsep=2pt]
\item 固定 $\mathscr T_{\Leftarrow}$ 时，最优压缩算子为 $\mathscr T_{\Rightarrow}^{\optstar}=(\mathscr T_{\Leftarrow}^{\star}\mathscr T_{\Leftarrow})^{-1}\mathscr T_{\Leftarrow}^{\star}$（标准实内积下矩阵形式为 $\bm T_{\Rightarrow}^{\optstar}=(\bm T_{\Leftarrow}^{\mathrm T}\bm T_{\Leftarrow})^{-1}\bm T_{\Leftarrow}^{\mathrm T}$），此时复合算子恰好还原正交投影 $\mathscr T_{\Leftarrow}\mathscr T_{\Rightarrow}^{\optstar}=\mathscr P_W$，使式~\eqref{cw:eq:manifold-orth-decomp} 次项归零；
\item 在全体 $n_d$ 维真子空间 $W<\mathcal S_X$ 中，前 $n_d$ 个主特征向量张成的子空间 $W^{\optstar}=\operatorname{span}\{u_1,\dots,u_{n_d}\}$（对应 $\bm P_{W^{\optstar}}=\sum_{j=1}^{n_d}\bm P_{u_j}$，其中 $\bm P_{u_j}:=u_ju_j^{\mathrm T}$）使重构误差达到全局极小值 $\sum_{j=n_d+1}^{n_D}\lambda_j$，且当 $\lambda_{n_d}>\lambda_{n_d+1}$ 时最优子空间唯一。
\end{enumerate}
\end{proposition}

\subsection{非线性参数化流形与推前先验}

当高维信号的内在几何呈显著弯曲形态（如手写数字图像随笔画形变构成的非线性曲面）时，平直子空间 $W$ 需要较高的维数才能压低残差特征值之和。为此，将压缩与还原算子推广为由参数 $\bm\theta\in\mathcal S_{\bm\Theta}$ 索引的神经网络映射族 $\mathscr T_{\Rightarrow}^{\bm\theta}:\mathcal S_X\to\mathcal S_Z$ 与 $\mathscr T_{\Leftarrow}^{\bm\theta}:\mathcal S_Z\to\mathcal S_X$，端到端重构映射记为 $\mathscr R^{\bm\theta}:=\mathscr T_{\Leftarrow}^{\bm\theta}\circ\mathscr T_{\Rightarrow}^{\bm\theta}$。还原算子的像集 $\mathcal M_{\bm\theta}:=\operatorname{im}(\mathscr T_{\Leftarrow}^{\bm\theta})$ 在环境空间 $\mathcal S_X$ 中给出由潜坐标 $z\in\mathcal S_Z$ 参数化的低维曲面；在任意潜坐标 $z_0\in\mathcal S_Z$ 的邻域内作一阶局部线性展开：
\begin{equation}\label{cw:eq:manifold-tangent}
\mathscr T_{\Leftarrow}^{\bm\theta}(z_0+\delta z)
\approx\mathscr T_{\Leftarrow}^{\bm\theta}(z_0)
+\frac{\mathrm d\mathscr T_{\Leftarrow}^{\bm\theta}}{\mathrm dz^{\mathrm T}}(z_0)\,\delta z,
\end{equation}
导数矩阵 $\frac{\mathrm d\mathscr T_{\Leftarrow}^{\bm\theta}}{\mathrm dz^{\mathrm T}}(z_0)\in\mathbb R^{n_D\times n_d}$ 的列空间即张成流形在该点处的局部切空间，当映射退化为常矩阵 $\bm T_{\Leftarrow}$ 时即回到命题~\ref{cw:prop:linear-manifold} 的平直情形。一般非线性参数下问题无闭式解，通过随机梯度优化联合极小化经验重构误差确定参数：
\begin{equation}\label{cw:eq:ml-reconstruction}
\bm\theta^{\optstar}\in\operatorname*{argmin}_{\bm\theta\in\mathcal S_{\bm\Theta}}\
\bigl\|(\mathscr I-\mathscr R^{\bm\theta})(x_{1..M})\bigr\|_{\hat{\mathbb P}_M}^2
=\operatorname*{argmin}_{\bm\theta\in\mathcal S_{\bm\Theta}}\
\frac1M\sum_{i=1}^M\bigl\|x_i-\mathscr T_{\Leftarrow}^{\bm\theta}\bigl(\mathscr T_{\Rightarrow}^{\bm\theta}(x_i)\bigr)\bigr\|^2.
\end{equation}
为避免过拟合或约束潜空间几何尺度，亦可在式~\eqref{cw:eq:ml-reconstruction} 中附加参数范数等正则项。下文固定训练所得参数 $\bm\theta^{\optstar}$，在不致混淆处将 $\mathscr T_{\Rightarrow}^{\bm\theta^{\optstar}},\mathscr T_{\Leftarrow}^{\bm\theta^{\optstar}}$ 简记为 $\mathscr T_{\Rightarrow},\mathscr T_{\Leftarrow}$。

确定压缩算子后，直接取训练样本经验测度 $\hat{\mathbb P}_{M;X}$ 在 $\mathscr T_{\Rightarrow}$ 下的推前测度作为潜状态的先验分布：
\begin{equation}\label{cw:eq:ml-prior}
\mathbb P_Z:=(\mathscr T_{\Rightarrow})_*\hat{\mathbb P}_{M;X}
=\frac1M\sum_{i=1}^M\delta_{\mathscr T_{\Rightarrow}(x_i)}.
\end{equation}
这一构造使潜空间初始粒子直接取自真实样本的压缩坐标，忠实保留了数据流形上的多峰聚集结构，无须强行假定潜变量服从标准正态分布。
